\section{Summary and conclusions}

 In this paper, we  have extended   the leading-logarithm
 analysis of lattice data for  parton pseudo-distributions
 and reduced pseudo-ITDs
 performed in Ref. \cite{Orginos:2017kos}.
 To this end, we incorporated
  recent  results for the reduced pseudo-ITDs at   the one-loop level   \cite{Radyushkin:2017lvu}
 (see also \cite{Ji:2017rah,Izubuchi:2018srq}).



 It was  found that the correction contains a  large term
 resulting  in essential numerical  changes  compared to  the LLA.
 The   large correction appears since
 effective distances involved in the most important diagrams
 are much smaller than the nominal distance $z_3$.
 This  leads to a  change
 (from $k_{\rm LLA} \approx 1$ to $k \approx 4$  in the case of our
 particular ITDs)  of the
 coefficient $k$ in  the
 rescaling relation $\mu = k/z_3$  that allows to (approximately)
 convert  the pseudo-PDFs ${\cal P} (x,z_3^2)$ into the $\overline{\rm MS}$
 PDFs $f(x,\mu^2)$.





While the  rescaling relation serves as an
instructive guide for quick estimates and  semi-quantitative analysis,
the $\overline{\rm MS}$ ITDs may be directly constructed applying   the
exact one-loop formula.  Using it, we have obtained
 the
   ITD ${\cal I} (\nu, \mu^2)$ at the \mbox{$\mu=1/a \approx 2.15$ GeV}
   $\overline{\rm MS}$ scale
         using the data in the  \mbox{$0\leq z_3 \leq 4a$}   region.

        We  found that ${\cal I} (\nu, \mu^2)$ at this scale is  close
         to the reduced pseudo-ITD ${\mathfrak M} (\nu,z_3^2)$  for $z_3\sim 4a$.
         Since  all the data in the $a\leq z_3 \leq 4a$ region do not differ much from
         the $z_3=4a$ ones (see Fig. \ref{M16}), the conversion of
         the ${\mathfrak M} (\nu,z_3^2)$ data into ${\cal I} (\nu, 1/a^2)$
         does not involve large changes, i.e.,
 the perturbative expansion for   $\overline{\rm MS}$  ITD ${\cal I} (\nu, \mu^2)$
 in terms of the reduced
   pseudo-ITDs ${\mathfrak M} (\nu,z_3^2)$    is  under control.

A  formal reason is that the   large  correction in this  case can  be
    absorbed into the $z_3^2$-dependent evolution term,
    with remaining corrections being small.

      Phenomenologically,
    the  PDF extracted in this way  follows the trend of  those
    given by the   global fits in the $x>0.1$ region,  but
    does not reproduce their  singular behavior in the $x<0.1$ region.
    The latter is usually related to the $x^{-0.5}$ pattern of the $\rho$-meson
    Regge trajectory. Since the $\rho$-meson is essentially a rather narrow
    resonance in the $\pi \pi$ system,
    one should not expect to accurately reproduce the $\rho$-meson
    properties  in a lattice simulation in which the pions
    are as heavy as 600 MeV.  Thus,  one may hope  that using simulations
    at physical pion mass would produce a better agreement with the global fits in the
    small-$x$ region.  This hope is supported by recent extractions  \cite{Alexandrou:2018pbm,Chen:2018xof}
    of $q_v (x)$ using  the quasi-PDF lattice simulations at physical pion mass.
